\documentclass[10pt,a4paper]{amsart}
\usepackage[margin=19mm]{geometry}
\usepackage{amsmath,amssymb,mathtools}
\usepackage[scr=rsfs,cal=euler]{mathalfa}
\usepackage{xfrac}
\usepackage[unicode,hidelinks,bookmarks=false]{hyperref}
\newcommand{\ie}{i.e.\ }
\newcommand{\cf}{cf.\ }
\newcommand{\resp}{respectively\ }
\newcommand{\Subsection}{\subsection}
\providecommand{\nobreakdash}{\nobreak\hbox{-}}
\setlength{\emergencystretch}{2em}
\multlinegap0pt
\usepackage{amssymb}
\newtheorem{definition}{Definition}
\newtheorem{lemma}{Lemma}
\newcommand{\F}{\mathcal{F}}
\newcommand{\M}{\mathcal{M}}
\newcommand{\Uu}{\mathcal{U}}
\title{Stability of Shifted Complexes via the Second-Moment Defect of the Up-Laplacian}
\author{Vinayak Gupta}
\date{Source: arXiv:2608.15358v1 (CC BY 4.0)}
\begin{document}
\maketitle

\noindent\emph{Source and licence.} Stability of Shifted Complexes via the Second-Moment Defect of the Up-Laplacian. Version arXiv:2608.15358v1 (CC BY 4.0), distributed by arXiv under the Creative Commons Attribution 4.0 International licence, http://creativecommons.org/licenses/by/4.0/, as recorded in the arXiv licence metadata for this exact version and shown on the abstract page https://arxiv.org/abs/2608.15358v1. Full licence notice: \texttt{licenses/arxiv-2608.15358-CC-BY-4.0.txt}.\par\smallskip

\noindent\textit{Editorial scope.} Counted unit: the article's lemma lem:Amonotone (source0.tex lines 331--336). The article's complete original proof of it is delivered with it (source0.tex lines 338--394, one \texttt{proof} environment of 57 lines). No mathematics is written, repaired or re-derived: the statement and the proof are the article's own lines, in the article's own notation, and no retained line is substituted or re-flowed.  Retained uncounted as the source-local context the proof needs: lines 287--298; lines 304--314.  Nothing was omitted from any retained span, so the package carries no omission record.  No retained line contains a citation, so the wrapper carries no bibliography and no bibliography entry.  No retained line contains a figure environment, an image-inclusion command or drawing code, so no image is delivered and the package declares no asset dependency.  Editorial formatting only, in addition: an AMS \texttt{amsart} preamble that restates the article's own declarations and macros used by the retained lines, namely the environments \texttt{definition}, \texttt{lemma} on separate counters, together with the standard packages those lines need.  The retained environments print in this document as Definition 1, Lemma 1 in order of appearance, whereas the article prints the same items with the numbering of its own full document.  The complete native source of the article as distributed in the arXiv e-print for this exact version is bundled unchanged as \texttt{sources/207738/source0.tex}.
\begin{definition}[Full compression]\label{def:compression}
Fix $i<j$ and let $\F\subseteq\binom{[n]}{r}$. Put
\[
\M=\M_{ij}(\F)=\{F\in\F:\ j\in F,\ i\notin F,\ \tau(F)\notin\F\},\qquad
\Uu=\F\setminus\M,
\]
where $\tau=\tau_{ij}$, and define the \emph{compression}
\[
C_{ij}\F=\Uu\cup\tau(\M).
\]
Thus a member is replaced by its shift exactly when that shift fails, and is left alone otherwise.
\end{definition}

\begin{lemma}\label{lem:basic}
Let $i<j$ and write $t=|\M|$. Then
\begin{enumerate}
\item[(a)] $t=N_{ij}(\F)-R_{ij}(\F)$;
\item[(b)] $\tau(\M)\cap\F=\emptyset$; in particular $\tau(\M)\cap\Uu=\emptyset$ and
$|C_{ij}\F|=|\F|$;
\item[(c)] writing $d'_v$ for the degrees in $C_{ij}\F$, one has $d'_i=d_i+t$, $d'_j=d_j-t$ and
$d'_v=d_v$ for $v\ne i,j$;
\item[(d)] $\tfrac12\bigl|\F\,\triangle\,C_{ij}\F\bigr|=t$.
\end{enumerate}
\end{lemma}

\begin{lemma}\label{lem:Amonotone}
For every $i<j$ and every $\F\subseteq\binom{[n]}{r}$,
\[
A\bigl(C_{ij}\F\bigr)\ \ge\ A(\F).
\]
\end{lemma}

\begin{proof}
Retain the notation $\M,\Uu,\tau$ of Definition \ref{def:compression} and write
$\F'=C_{ij}\F=\Uu\cup\tau(\M)$. We construct an injection $\Phi$ from the set of adjacent pairs of
$\F$ into the set of adjacent pairs of $\F'$, where a pair is \emph{adjacent} when its two members
are distinct and meet in $r-1$ points. Let $\{F,G\}$ be an adjacent pair of $\F$. Since
$\F=\Uu\sqcup\M$, there are three cases.

\smallskip
\noindent\textit{Case 1: $F,G\in\Uu$.} Both members survive the compression, so
$\Phi(\{F,G\})=\{F,G\}$ is an adjacent pair of $\F'$ consisting of two old members.

\smallskip
\noindent\textit{Case 2: $F,G\in\M$.} Both are replaced, and
$|\tau(F)\cap\tau(G)|=|F\cap G|=r-1$ because $\tau$ is a bijection of $[n]$. Moreover
$\tau(F)\ne\tau(G)$. So $\Phi(\{F,G\})=\{\tau(F),\tau(G)\}$ is an adjacent pair of $\F'$ consisting
of two new members.

\smallskip
\noindent\textit{Case 3: $F\in\M$ and $G\in\Uu$.} Write $F=X\cup\{j\}$ with $i,j\notin X$. We
distinguish three sub-cases according to how $G$ meets $\{i,j\}$.

\smallskip
\noindent\textit{(3a): $G$ contains both of $i,j$ or neither of them.} Then $\tau(G)=G$, so
\[
|\tau(F)\cap G|=|\tau(F)\cap\tau(G)|=|F\cap G|=r-1 .
\]
Also $\tau(F)\ne G$, since $\tau(F)\notin\F$ while $G\in\F$. As $G\in\Uu$ survives, we may set
$\Phi(\{F,G\})=\{\tau(F),G\}$, an adjacent pair of $\F'$ with exactly one new member.

\smallskip
\noindent\textit{(3b): $j\in G$ and $i\notin G$.} Since $G\in\Uu$ was not moved, the only possible
reason is $\tau(G)\in\F$. The set $\tau(G)$ contains $i$ and not $j$, hence $\tau(G)\notin\M$ and
therefore $\tau(G)\in\Uu$; in particular $\tau(G)$ is an old member of $\F'$. As in Case 2,
$|\tau(F)\cap\tau(G)|=r-1$ and $\tau(F)\ne\tau(G)$, so we may set
$\Phi(\{F,G\})=\{\tau(F),\tau(G)\}$, again an adjacent pair of $\F'$ with exactly one new member.

\smallskip
\noindent\textit{(3c): $i\in G$ and $j\notin G$.} We claim this cannot occur. Write $G=Y\cup\{i\}$
with $i,j\notin Y$, so that $|X|=|Y|=r-1$. Since $j\notin G$ and $i\notin F$, we get
$F\cap G=X\cap Y$, and $|F\cap G|=r-1$ forces $X=Y$. Hence $G=X\cup\{i\}=\tau(F)$. But $F\in\M$
means $\tau(F)\notin\F$, contradicting $G\in\F$.

\smallskip
This defines $\Phi$ on all adjacent pairs of $\F$, and every image is an adjacent pair of $\F'$. It
remains to prove that $\Phi$ is injective. By Lemma \ref{lem:basic}(b) the old and new members of
$\F'$ are distinguishable, and the images produced in Cases 1, 2 and 3 contain respectively zero,
two and exactly one new member. Hence no two pairs from different cases have the same image, and it
suffices to check injectivity within each case.

In Case 1 the map is the identity. In Case 2 it is induced by the bijection $\tau$, hence injective.
In Case 3 an image has the form $\{\tau(F),Z\}$ with $\tau(F)$ new and $Z\in\Uu$ old; since $\tau$ is
injective, the new entry determines $F$. It determines the old entry too: in sub-case (3a) we have
$Z=G$, which contains both of $i,j$ or neither, whereas in sub-case (3b) we have $Z=\tau(G)$, which
contains $i$ but not $j$. These two possibilities are mutually exclusive, so $Z$ determines which
sub-case occurred, and then $G=Z$ or $G=\tau(Z)$ accordingly. Thus $\Phi$ is injective on Case 3 as
well, and $A(\F')\ge A(\F)$.
\end{proof}

\end{document}
